\section{Discussion}

The analysis presented demonstrates the detection of GW150914 using open LIGO data. The measured time delay between the H1 and L1 detectors, approximately 7 ms, is consistent with the light-travel time across the Earth, confirming the coherence of the signal across detectors. Time–frequency analysis using the Q-transform reveals the characteristic chirp of a binary black hole merger, with increasing frequency and amplitude, which is consistent with theoretical predictions. 

Matched filtering with an optimized waveform template yields a peak reweighted SNR of approximately 13.3, indicating a strong detection well above noise fluctuations. The reweighted SNR incorporates the results of the $\chi^2$ signal consistency test, which measures deviations from the template and helps mitigate the influence of transient noise. The best-matching template corresponds to component masses $M_1 = 35.5 M_\odot$ and $M_2 = 30.5 M_\odot$, resulting in a total mass of $66 M_\odot$, in agreement with published results \cite{Abbott2016Discovery}. The reduced $\chi^2$ value of 2.10 indicates an acceptable agreement between the data and the template, confirming the astrophysical origin of the signal with a detection significance of roughly 13$\sigma$.

Although the $\chi^2$ test helps suppress spurious noise contributions, residual short-duration noise transients, or glitches, may still affect the data. Future analysis could focus on additional noise mitigation techniques, such as vetoing coincident triggers in auxiliary channels or applying machine learning-based classifiers to identify and remove glitches before matched filtering. Such improvements would increase the robustness of gravitational-wave detections and enhance confidence in low-SNR events.

Overall, the combination of time–frequency and matched-filter analyses provides a reliable confirmation of GW150914 and demonstrates the effectiveness of open data analysis tools for gravitational-wave astrophysics.
